\section{Conclusiones y trabajo futuro}

Este informe cerr\'o la bater\'ia m\'inima del laboratorio sin bono con seis experimentos reproducibles y una l\'inea principal cl\'asica sobre EEG monocanal. Los resultados permiten tres conclusiones principales. Primero, dentro de dataset, la combinaci\'on de ventanas de 30 s, \feature{eeg\_*} espectrales/de complejidad y ensambles basados en \'arboles ofrece baselines s\'olidos. En Sleep-EDF, XGBoost fue el mejor por \emph{accuracy} y \emph{kappa}, mientras que Random Forest fue el mejor por \emph{macro-F1}. En MIT-BIH multitarget, XGBoost fue el mejor compromiso medio entre apnea y staging.

Segundo, el principal problema no es la optimizaci\'on interna sino la \textbf{generalizaci\'on}. Todos los experimentos \emph{cross-dataset} mostraron degradaciones fuertes. En apnea, incluso aparecieron soluciones con sensibilidad nula y especificidad perfecta, una combinaci\'on inaceptable para despliegue real aunque la \emph{accuracy} parezca alta. Por ello, el hallazgo m\'as importante del trabajo es que los baselines cl\'asicos son \'utiles como referencia reproducible, pero no bastan por s\'i solos para sostener transferencia entre cohortes.

Tercero, el an\'alisis m\'etrico confirma que no basta con mirar \emph{accuracy}. La baja F1 de N1, la sensibilidad inestable de apnea, la superposici\'on parcial de intervalos de confianza y la discrepancia entre media de folds y AUC agregada muestran que las conclusiones deben ser matizadas y apoyarse en varias evidencias a la vez.

Como trabajo futuro se proponen tres l\'ineas directas: incorporar un m\'odulo expl\'icito de manejo de artefactos, extender el an\'alisis a severidad cl\'inica por paciente cuando el \emph{ground truth} sea consistente, e implementar una estrategia de adaptaci\'on de dominio para EEG monocanal. Esa tercera l\'inea es la m\'as prometedora, porque ataca precisamente el cuello de botella que la bater\'ia experimental final dej\'o al descubierto.
